% Folder 156-8, batch 05 — archivists' pages 81–100 (author's pp. 73–92).
% Pass: Opus 5.5 (claude-opus-5-5), 2026-10-02 — first pass, unchecked against the pages by a human.
%
% Page numbers follow the project's convention (page 1 of batch-05.pdf = 81,
% as the tiles are named). The pencilled numbers bottom-left on the leaves
% read one higher (82 on the first sheet, 101 on the last); the offset was
% left as is so that \page{} matches the facsimile the site turns.
% Notation fixed for this batch: supp°, cosupp°, Omb, Omb° as \mathrm{…};
% |∘| (disjonction) as \mathrel{|{\circ}|}; its barred negation as
% \mathrel{\overline{|\circ|}}; ≪ with a ring above as \overset{\circ}{\ll};
% the compatibility sign (a crossed ≶) as \not\lessgtr; 𝔉 for the set of
% figures, \mathcal{M} for the magasin, Drap(L) for the flags of L.
\documentclass[11pt,a4paper]{article}
\input{../preamble/grothendieck}

\folder{156-8}
\batch{5}
\pages{81}{100}
\foldertitle{[Chapitre] VIII. Analysis situs (quatrième mouture) : notes manuscrites (26/06-04/07/1986).}
\dating{1986}
\watermark{Édition de démonstration}

\begin{document}

\section{Supports intérieurs et fonctions support pleinement fidèles}

\page{81}
\note{p. 73 de l'auteur. La page reprend un argument commencé avant ce lot : les notations $\mathcal{M}$, $|X|^\circ$, $\mathrm{supp}^\circ$, $\varphi$ y sont déjà en usage.}

\struck{Soit. Supposons}

\struck{Soit $F$ \add{$\subset\mathcal{M}$} une figure de $\mathcal{M}$. La chose essentielle, ici, c'est que}
\[
\text{\struck{$X,Y\in F \Rightarrow X \not\lessgtr Y$}}
\]
\struck{i.e. que $F$ est contenue dans une figure, (et une \add{\uncertain{partielle}} \uncertain{telle} est \ill{} une figure vraiment, i.e. fermée pour $\le$).}
\note{ce premier paragraphe est barré de trois traits obliques. Le signe de compatibilité, un $\lessgtr$ barré, est lu $\not\lessgtr$ dans tout le lot.}

Je vais m'intéresser d'abord, plus particulièrement, à
\[
\mathrm{supp}^\circ X = A \subset \mathcal{M} \qquad (X\in\mathcal{M})
\]
et sa relation à
\[
|X|^\circ \subset P .
\]
On a $X \in \mathrm{supp}^\circ X$ donc
\[
|X|^\circ \subset |\mathrm{supp}^\circ X|^\circ .
\]
J'aimerais que l'\add{inclusion} inverse soit vraie, i.e. que
\[
|X|^\circ = |\mathrm{supp}^\circ X|^\circ
\]
ou encore que
\[
\forall Y \in \mathrm{supp}^\circ X, \text{ on ait } |Y|^\circ \subset |X|^\circ
\]
(l'implication inverse étant claire). Il revient au même de dire que \struck{$Y\notin\mathrm{supp}$}
\[
\begin{gathered}
|Y|^\circ \not\subset |X|^\circ \Longrightarrow Y \notin \mathrm{supp}^\circ X, \text{ i.e.}\\
\exists Z\in\mathcal{M} \text{ avec } Z \mathrel{|{\circ}|} X,\; Z \mathrel{\overline{|\circ|}} Y,
\end{gathered}
\]
i.e. $|Y|^\circ \cap \complement |X|^\circ \neq \emptyset$ à gauche, et à droite
\[
\begin{cases}
|Z|^\circ \cap |X|^\circ = \emptyset\\
|Z|^\circ \cap |Y|^\circ \neq \emptyset .
\end{cases}
\]

\page{82}
\note{p. 74 de l'auteur.}

C'est \uncertain{clair} \struck{\ill{}} dans le cas où
\[
|Y|^\circ \cap \partial X \neq \emptyset ,
\]
\struck{i.e. \uncertain{quand} $\exists Z \le X$, \ill{}} ce qui signifie, puisque
\[
\partial X = \bigcup_{Z<X} |Z|^\circ ,
\]
que
\[
\exists Z < X, \quad |Y|^\circ \cap |Z|^\circ \neq \emptyset ,
\]
on prend donc ce $Z$-là \uncertain{disons}, on aurait donc
\[
|Y|^\circ \cap \complement |X| \neq \emptyset
\]
\struck{\uncertain{pour}} puisque $\complement |X|^\circ = |\partial X| \sqcup \complement |X|$, et on \ill{} \ill{} fait l'hypothèse \ill{}

\textbf{Hyp.} $\forall X\in\mathcal{M}$, \struck{on} et tout $x \in P \setminus |X|$,
\[
\exists Z \in \mathcal{M} \text{ \struck{$\complement|X|$} avec } x\in|Z|, \quad Z \parallel X \quad (\text{i.e. } |Z|\cap|X| = \emptyset) .
\]
\marginal{NB quitte à remplacer $Z$ par un de ses \ill{}, on peut supposer \ill{} $x \in |Z|^\circ$.}
\marginal{C'est \uncertain{tout fait}, il est plus raisonnable d'écrire $Z \mathrel{|{\circ}|} X$ (pour \uncertain{inclure} p.ex. les \uncertain{magasins} \uncertain{fidèles}, et notamment \ill{} \ill{} les \ill{}).}

Pour ce qui m'occupe à présent, il suffirait \struck{\ill{}} d'avoir $Z \mathrel{|{\circ}|} X$, i.e. $|Z|^\circ \cap |X|^\circ = \emptyset$. Cette hypothèse équivaut (pour $X$ donné) à
\[
P \setminus |X| = \bigcup_{Z\in\mathrm{cosupp}^\circ(X)} |Z| \qquad \text{ou}
\]
\[
\Bigl(= \bigcup_{Z\in\mathrm{cosupp}^\circ(X)} |Z|^\circ\Bigr)
\]
\note{l'exposant de « cosupp » sous les deux réunions est un petit signe peu lisible, lu ici $^\circ$.}

\page{83}
\note{p. 75 de l'auteur.}

\struck{En fait, pour le moment il me suffirait d'avoir} Ceci implique
\[
P \setminus |X|^\circ = \bigcup_{Z\in\mathrm{cosupp}^\circ X} |Z|^\circ
\]
\marginal{NB Avec cette seule relation, il n'est pas exclu que parmi les $Z$ qui interviennent, il y en ait avec $|Z|\cap|X|^\circ \neq \emptyset$ !}
et ceci me suffit pour l'instant. On voit que ça assure déjà les deux relations
\[
\begin{aligned}
\varphi(\mathrm{supp}^\circ X) &\overset{\text{déf}}{=} |\mathrm{supp}^\circ X|^\circ = |X|^\circ\\
\varphi(\mathrm{cosupp}^\circ X) &\overset{\text{déf}}{=} |\mathrm{cosupp}^\circ X|^\circ = P\setminus |X|^\circ .
\end{aligned}
\]
\struck{qui impliquent} Je dis qu'on a aussi \struck{$\mathrm{supp}\,X$}
\[
\begin{aligned}
\varphi(\mathrm{supp}\,X) &= |X|\\
\varphi(\mathrm{cosupp}\,X) &= P\setminus |X|
\end{aligned}
\]
si \struck{\ill{}} par définition
\[
\begin{aligned}
\mathrm{supp}\,X &= \mathrm{supp}^\circ(\widetilde{X})\\
\mathrm{cosupp}\,X &= \mathrm{cosupp}^\circ(\widetilde{X})
\end{aligned}
\]
Comme pour $Y\in\widetilde{X}$, on a $Y\in\mathrm{supp}\,X$, d'où \struck{$\varphi(Y)\subset$} $|Y|^\circ \subset \varphi(\mathrm{supp}\,X)$ \uncertain{i.e. tant} que $\varphi(\mathrm{supp}\,X)$ contient les $|Y|^\circ$, donc leur réunion $|X|$ :
\[
|X| \subset \varphi(\mathrm{supp}\,X) .
\]
Pour l'implication inverse, il faut voir que
\[
Y \in \mathrm{supp}\,X \;(= \mathrm{supp}^\circ(\widetilde{X})) \Longrightarrow |Y|^\circ \subset |X|
\]
or \uncertain{comme tantôt}
\[
|Y|^\circ \not\subset |\overline{X}| \Longrightarrow Y \notin \mathrm{supp}^\circ(\widetilde{X}) \quad \text{i.e. } \exists Z
\]
\[
\text{i.e. } |Y|^\circ \cap (P \setminus \overline{X}) \neq \emptyset
\qquad
\text{avec } X \mathrel{|{\circ}|} Z,\; Y \mathrel{\overline{|\circ|}} Z ,
\]
\[
\text{i.e. } |X| \cap |Z|^\circ \neq \emptyset,\; |Y|^\circ \cap |Z|^\circ .
\]
\note{sic : on attend $|X|\cap|Z|^\circ = \emptyset$, $|Y|^\circ\cap|Z|^\circ \neq \emptyset$ ; la dernière ligne, serrée au bas du feuillet, est transcrite telle qu'elle se lit, la seconde relation sans son second membre. $X$ y est écrit tantôt $\widetilde{X}$, tantôt $\overline{X}$.}

\page{84}
\note{p. 76 de l'auteur.}

\add{dont} \uncertain{c'est} ici qu'on aura besoin de l'hypothèse \add{*} plus haut, dans toute sa force.

Mais je soupçonne que ce n'est pas \ill{} \ill{} — comme un exemple, \ill{}

Soit $Y < X$, et considérons
\[
\mathrm{supp}^\circ(\{Y,X\}) \quad\text{et}\quad \varphi(\mathrm{supp}^\circ(\{Y,X\}))
\]
on aura donc
\[
\text{\struck{\ill{}}}\; |X|^\circ \cup |Y|^\circ \subset \varphi(\mathrm{supp}^\circ(\{Y,X\}))
\]
aura-t-on l'inverse ? \struck{Je n'en} \add{Mais oui} — et je vois \add{à présent} quelle est l'hypothèse générale qui \struck{\ill{}} \add{\uncertain{assure}} qu'on a une « bonne compatibilité » pour des « supports \add{(intérieurs)} constructibles ».
\marginal{(* C'est \ill{} \ill{} jusqu'ici \ill{} fait \ill{} $P$ \ill{} ici, si deux \ill{} $Z$ \ill{} $|Z| \subset P\setminus|X|$ i.e. $|Z|^\circ\cap|X| = \emptyset$, on \ill{} \ill{})}

\textbf{Définition} Soit $\mathfrak{F}$ un \ill{} associé au magasin $\mathcal{M}$. Un élément \add{$A$} de $\Sigma_{\mathcal{M}}$ on dit un support (intérieur) \textit{constructible}, s'il existe une \struck{figure $F\in\mathfrak{F}$ et une} \add{(i.e. $\Phi\subset F$, avec $F\in\mathfrak{F}$)} partie $\Phi$ de $\mathcal{M}$, majorée dans $\mathfrak{F}$, telle que l'on ait
\[
A = \text{\struck{\ill{}}}\; \mathrm{supp}^\circ(\Phi) .
\]
On dit qu'une fonction support \uncertain{immédiatement} extérieure [fidèle]
\[
X \longmapsto |X|, \quad \mathcal{M} \longrightarrow \mathfrak{P}(P)
\]
est \textit{pleinement fidèle} (pour $\mathfrak{F}$), si elle

\page{85}
\note{p. 77 de l'auteur.}

satisfait la propriété suivante
\[
\forall F \in \mathfrak{F}, \text{ on a } \varphi(\mathrm{cosupp}^\circ(F)) = X \setminus |F|
\]
i.e.
\[
X\setminus|F| = \bigcup_{\substack{Z\in\mathcal{M}\\ \text{t.q. } Z \mathrel{|{\circ}|} F}} |Z|^\circ
\qquad [\text{i.e. } |Z|^\circ \subset X\setminus|F|]
\]
i.e.
\[
\boxed{\forall x \in X - |F|, \;\exists Z \text{ avec } x \in |Z|^\circ \subset X\setminus|F|}
\]
\marginal{? plutôt $P$}
\note{les deux $X$ de « $X\setminus|F|$ » sont entourés et reliés à la note marginale « plutôt $P$ ».}
\marginal{Sous cette forme, la propriété \uncertain{implique} la fidélité.}

Ceci posé, on trouve

\textbf{Proposition} Soit $\Phi \subset F$. Alors
\[
\varphi(\mathrm{supp}^\circ\Phi) = |\Phi|^\circ \Bigl(\overset{\text{déf}}{=} \bigcup_{X\in\Phi} |X|^\circ\Bigr) \qquad (*)
\]

\textbf{Dém} Comme \struck{$X\in$} $\Phi \subset \mathrm{supp}^\circ(\Phi)$, on a \struck{$\varphi(\Phi)=$}
\[
\varphi(\Phi) \overset{\text{déf}}{=} |\Phi|^\circ \subset \varphi(\mathrm{supp}^\circ(\Phi)) ,
\]
i.e. on a $\subset$, \struck{Soit $Y\in\mathrm{supp}^\circ\Phi$}, i.e. Prouvons $\supset$.
\[
\text{\uncertain{Soit}}\quad Y \in \mathrm{supp}^\circ(\Phi) \Longrightarrow |Y|^\circ \subset |\Phi|^\circ \quad \text{i.e.}
\]
\[
\begin{gathered}
(*)\quad |Y|^\circ \not\subset |\Phi|^\circ \Longrightarrow Y \notin \mathrm{supp}^\circ(\Phi), \text{ i.e. } \exists Z \ldots\\
Z \mathrel{|{\circ}|} \Phi, \text{ et } Z \mathrel{\overline{|\circ|}} Y, \text{ i.e.}\\
|Z|^\circ \cap |\Phi|^\circ = \emptyset,\; |Z|^\circ\cap|Y|^\circ \neq \emptyset .
\end{gathered}
\]
i.e. $|Y|^\circ \cap \complement|\Phi|^\circ \neq \emptyset$.
\marginal{* explicitons d'abord la relation $\Phi \mapsto |\Phi|^\circ$, \ill{} en termes de \ill{} figures \ill{} \ill{}}

Distinguons deux cas

(i) $|Y|^\circ \cap (|F| \setminus |\Phi|^\circ) \neq \emptyset$, \struck{si $x$ est dans le}
\[
|F|\setminus|\Phi|^\circ = |F\setminus\Phi|^\circ ,
\]
$\exists Z \in F\setminus\Phi$, tel que $|Y|^\circ \cap |Z|^\circ \neq \emptyset$. OK

(ii) $|Y|^\circ \cap (P\setminus|F|) \neq \emptyset$, alors on applique l'hypothèse de pleine fidélité à $F$.

\page{86}
\note{p. 78 de l'auteur.}

\textbf{Corollaire 1} L'application
\[
\Phi \longmapsto \mathrm{supp}^\circ(\Phi) : \mathfrak{P}(F) \longrightarrow \Sigma_{\mathcal{M}}
\]
commute aux Inf et Sup
\[
\begin{aligned}
\mathrm{supp}^\circ\Bigl(\bigcap_i \Phi_i\Bigr) &= \bigcap_i \mathrm{supp}^\circ(\Phi_i) && I \neq \emptyset\\
\mathrm{supp}^\circ\Bigl(\bigcup_i \Phi_i\Bigr) &= \operatorname{Sup}_i(\mathrm{supp}^\circ(\Phi_i)) && \text{sup dans } \Sigma_{\mathcal{M}}
\end{aligned}
\]
et de plus \uncertain{on a}, si $\Phi \subset \Psi$
\[
\mathrm{supp}^\circ(\Psi\setminus\Phi) = \mathrm{supp}^\circ\Psi \cap \complement\,\mathrm{supp}^\circ\Phi
\]

\textbf{Corollaire 2} \struck{Dans} \add{Considérons} l'image $\Sigma_{\mathcal{M},F}$ de $\mathfrak{P}(F)$ (\uncertain{ensemble} des supports constructibles \uncertain{subordonnés} à $F$). Cette image est stable dans $\Sigma_{\mathcal{M}}$, par inf (familles non vides) et sup, et complémentaire relatif, i.e. opération $A \cap \complement B$. De plus, l'application \struck{$A \mapsto |A|^\circ$}
\[
\begin{gathered}
A \longmapsto |A|^\circ = \bigcup_{X\in A} |X|^\circ\\
\Sigma_{\mathcal{M},F} \longrightarrow \mathfrak{P}(|F|) \subset \mathfrak{P}(P)
\end{gathered}
\]
commute à ces opérations.

\textbf{Terminologie} \uncertain{Une famille} de supports constructibles \add{(\ill{} $F$)} est dite « admissible », s'il existe $F\in\mathfrak{F}$, tel qu'ils soient tous subordonnés à $F$. \struck{l'image de $\mathfrak{P}(F)$ dans $\Sigma_{\mathcal{M}}$.}

\textbf{Scholie} Pour les familles admissibles de supports constructibles, les opérations Sup, Inf, complémentaire relatif — ont toutes

\page{87}
\note{p. 79 de l'auteur.}

les bonnes propriétés \struck{\uncertain{habituelles}} ensemblistes habituelles, notamment
\[
\begin{aligned}
A \cap \operatorname{Sup}_i B_i &= \operatorname{Sup}_i (A\cap B_i)\\
A \setminus (A\setminus B) &= \text{\struck{$A\cap$}}\, B \qquad (\text{si } B\subset A)
\end{aligned}
\]
De plus \struck{$A\cap B$}
\[
A\cap B = \emptyset \Longleftrightarrow A \mathrel{|{\circ}|} B ;
\]
(et pas seulement $\Leftarrow$) Toute pathologie disparaît.

\textbf{Question} Soit $\mathcal{M}$ un magasin quelconque, \struck{\ill{}} supposons l'existence d'une fonction support extérieure pleinement fidèle. Le scholie sur les familles admissibles de \uncertain{supports} constructibles, \struck{\ill{}} \add{i.e.} \ill{} des \uncertain{axiomes} \ill{}, \uncertain{sont-ils} nécessairement valables ?
\marginal{Non}
Aussi : \struck{existe} des fonctions support fidèles \uncertain{existent-elles} qui ne soient pleinement fidèles ??

\textbf{Axiome} \struck{Ax} supp 1 (« axiome des supports »). Soient : $F\in\mathfrak{F}$, $\Phi\subset F$, $\Phi' = F\setminus\Phi$. Alors :
\[
\mathrm{supp}^\circ(F) \cap \mathrm{cosupp}^\circ\Phi = \mathrm{supp}^\circ(\Phi')
\]
en d'autres termes : $\forall Z\in\mathcal{M}$
\[
\text{\struck{si}}\; Z \in \mathrm{supp}^\circ(F),\; \text{\struck{et}}\; Z \parallel \Phi \Longrightarrow Z \in \mathrm{supp}^\circ(\Phi')
\]
\note{après « $\mathrm{supp}^\circ(F)\cap$ », un mot noirci et rayé précède « cosupp° ».}

NB Cet axiome ne fait intervenir que $\le$, $|{\circ}|$ \ill{} et non pas $\overset{\circ}{\ll}$.

\page{88}
\note{p. 80 de l'auteur.}

C'est le premier axiome de nature délicate, sur lequel nous tombions. Il est automatiquement satisfait par les magasins fidèles (car $\mathrm{Omb}$ y est une fonction support ext.\ pleinement fidèle), plus généralement chaque fois qu'il y a une fonction support extérieure pleinement fidèle. Le premier cas intéressant où il n'est pas sûr qu'on ait une telle fonction support, est celui des magasins associés à un \struck{\ill{}} ensemble ordonné, qu'il serait temps d'introduire ici, et d'examiner sous l'optique des \uncertain{figures} et des supports.

Je voudrais donc à présent

1°) regarder le cas des magasins \struck{associés} définis par un ens.\ ordonné

2°) Examiner le cas des magasins définis par $\le$, $|{\circ}|$ seulement ($\overset{\circ}{\ll}$ étant défini via $\le$, $|{\circ}|$, \struck{\ill{}} \add{cette description} permettant de donner un sens à la relation $x^\circ = \mathrm{supp}^\circ(x,\ldots)$, donc $y^\circ \subset x^\circ$ (en définissant $\overset{\circ}{\ll}$ …) en définissant en conséquence

3°) Voir si dans ces magasins « disjonctifs », on \uncertain{n'arrivait} pas, par des exemples évidents, à ce que l'axiome des supports serait en défaut…*
\marginal{* Oui !}

\section{Magasin associé à un ensemble ordonné}

\page{89}
\note{p. 81 de l'auteur.}

\textbf{1 juillet} Magasin associé à un ensemble ordonné $L$.

A) \struck{\ill{}} Supposons d'abord $L$ \textit{totalement ordonné}, et \textit{divisible} i.e.
\[
\forall a,b\in L, \quad a<b \Longrightarrow \;]a,b[\; \neq \emptyset
\qquad \Bigl(]a,b[ \;\overset{\text{déf}}{=} \{x\in L \mid a<x<b\}\Bigr)
\]
Dans ce cas, on va définir un \textit{magasin unidimensionnel}, ayant $L$ comme ens.\ des points (ou lieux), et avec
\[
\mathcal{M} \subset \text{\struck{$\mathfrak{P}(\mathfrak{P}(L))$}}\; \mathrm{Fig}(L)
\]
défini ainsi :
\[
\begin{cases}
\mathcal{M} = \mathcal{M}_0 \sqcup \mathcal{M}_1, \quad \text{avec}\\
\mathcal{M}_0 = \{P_x \mid x \in L\}, \quad \text{où } P_x = \{\{x\}\}\\
\mathcal{M}_1 = \{S_{a,b} \mid a,b\in L,\; a<b\}, \quad S_{ab} = \{[a,b],\{a\},\{b\}\}
\end{cases}
\]
\marginal{NB \ill{} \ill{} écrire \ill{} Sans \ill{} $a<b$ \ill{} de divisibilité \ill{} $S_{a,b}$ \ill{} \ill{} figures, i.e. $|S_{a,b}|^\circ \neq \emptyset$ \ill{}}
Il est tautologique que les conditions Mag 1, 2 sont vérifiées. ($\mathcal{M}$ contient les $\{\{x\}\}$ i.e.\ les $X_x$, et $\mathcal{M}$ stable par passage à une sous-\struck{multi} multistructure.)

Explicitons les relations $\le$, $\overset{\circ}{\ll}$, $|{\circ}|$

1°) $X\le Y \Longleftrightarrow X = Y$ ou $X = P_x$, $Y = S_{a,b}$, avec $x\in\{a,b\}$, i.e. $x = a$ ou $x = b$.

\struck{2°) $X \overset{\circ}{\ll} Y$ ssi $X = Y$ ou $X = P_x$, $Y$} en d'autres termes, on a
\[
\begin{cases}
\partial P_x = \emptyset\\
\partial S_{a,b} = \{P_a, P_b\}
\end{cases}
\]

2°) $X \overset{\circ}{\ll} Y \Longleftrightarrow$ \struck{$X = Y$ ou $X = P_a$}
\[
\begin{cases}
\text{a) si } X = P_x,\, Y = P_y : & x = y \quad \text{i.e. } P_x = P_y\\
\text{b) si } X = P_x,\, Y = S_{ab} : & a<x<b \quad [\text{i.e. } x\in|X|^\circ]\\
\text{c) si } X = S_{a,b},\, Y = S_{a',b'} : & a'\le a<b\le b'
\end{cases}
\]
\[
[\text{i.e. } |X|\subset|Y|,\; |X|^\circ\subset|Y|^\circ]
\]
\note{en b), on attend $x\in|Y|^\circ$. En c), quelques mots rayés sous « $a'\le a<b\le b'$ » ne se lisent pas.}

On avait glané* que les $P_x$ sont les éléments minimaux pour $\le$, et les éléments minimaux pour $\overset{\circ}{\ll}$ (grâce à divisibilité), ce sont donc aussi les él.\ minimaux pour $\overset{\circ}{\ll}$ = lieux.
\marginal{NB Je vais \uncertain{maintenant} utiliser la notation $X \mathrel{\Delta} Y$ pour $X<Y$ et $X \mathrel{\underline{\Delta}} Y$ pour $X\le Y$ (notations d'\uncertain{incidence} \ill{} visuellement \ill{}) pour $X<Y$ à la notation « positionnelle » $X \le_{\mathrm{pos}} Y$.}
\marginal{* c'est d'ailleurs \ill{} un \ill{} des \ill{} \uncertain{Magasins} 1, 2. Voir p. 97, 98.}
\note{p. 97, 98 : pagination de l'auteur, au-delà de ce lot.}

\page{90}
\note{p. 82 de l'auteur.}

3°)
\[
X \mathrel{|{\circ}|} Y \Longleftrightarrow
\begin{cases}
X = P_x,\, Y = P_y : & [x\neq y \text{ i.e.}]\; x<y \text{ ou } y<x\\
X = P_x,\, Y = S_{a,b}\;(\text{ou l'inverse}) : & x\le a \text{ ou } b \le x\\
X = S_{a,b},\, Y = S_{a',b'} : & a<b\le a'<b' \text{ ou } a'<b'\le a<b
\end{cases}
\]
\note{au-dessus de l'accolade, deux signes peu lisibles (« $\partial L$ » ?). Dans le deuxième cas, un $\ge$ est surchargé en $\le$.}

On a exprimé simplement $|X|^\circ \cap |Y|^\circ$ \struck{\ill{}} $= \emptyset$, en utilisant la divisibilité.

Dans le contexte ensembliste, il n'y a \uncertain{même} plus d'axiomes à vérifier, par \ill{} plus d'axiome \struck{\ill{}} de la fonction support, grâce à l'existence de la fonction \add{pour} support extérieure fidèle $X \mapsto |X|$ (comme pour tous magasins ensemblistes, \uncertain{simplement} partiel, voire, seulement fidèle — on prend cette fois la fonction support fidèle $X \mapsto \mathrm{Omb}\,X$).

B) Supposons maintenant la relation d'ordre quelconque. On prendra encore
\[
\begin{cases}
\mathcal{M} \subset \mathfrak{P}(\mathfrak{P}(L))\\
\mathcal{M} = \mathcal{M}_0 \sqcup \mathcal{M}_1\\
\mathcal{M}_0 = \{P_x \mid x\in L\}, \quad P_x = \{\{x\}\}\\
\mathcal{M}_1 = \{S_{a,b} \mid a,b\in L,\; a<b\}, \quad S_{a,b} = \{[a,b],\{a\},\{b\}\}
\end{cases}
\]
\marginal{il \struck{\ill{}} serait plus simple de prendre $\mathcal{M}_0 = L$, $\mathcal{M}_1 = \{(a,b)\in L\times L \mid a<b\}$}
mais on notera qu'il ne s'agit pas ici de figures (élémentaires), faute de savoir si
\[
|S_{a,b}|^\circ = [a,b]\setminus\{a,b\} = \;]a,b[
\]
est non vide. On \textit{définira} alors les relations $\le$, $\overset{\circ}{\ll}$, $|{\circ}|$ par les formules explicitées plus haut, en faisant \uncertain{remarquer} cette fois que
\[
S_{a,b} \overset{\circ}{\ll} S_{a',b'} \Longleftrightarrow a'\le a<b\le b' \Longleftrightarrow |S_{a,b}|^* \subset |S_{a',b'}|^*
\]
\[
\Longrightarrow |S_{a,b}|^\circ \subset |S_{a',b'}|^\circ
\]
[mais l'implication inverse pas nécessairement valable
\note{au-dessus du second $\Longleftrightarrow$ est écrit un $\Rightarrow$ barré ; sous la formule, une ligne raturée illisible ; les exposants $^*$ sont lus tels quels.}
\marginal{NB Avec la présentation adoptée, on trouve des « pré-figures » unidimensionnelles dans $L$ (où \uncertain{peut} avoir $|X|^\circ = \emptyset$) (on y voit que $\overset{\circ}{\ll}$ est \uncertain{réduit} à $\le$, \ill{} \ill{} par les relations ensemblistes. Mais pour $|{\circ}|$) \ill{} dans le cas tot.\ ordonné.}

\struck{$P_x \mathrel{|{\circ}|} P_y \Longleftrightarrow x<y$ ou}

\page{91}
\note{p. 83 de l'auteur.}

\[
P_x \mathrel{|{\circ}|} P_y \Longleftrightarrow x<y \text{ ou } y<x \text{ i.e. } \{x,y\}\in\mathrm{Drap}_2(L) \Longrightarrow x\neq y
\]
$[\Longleftrightarrow |P_x|^\circ \cap |P_y|^\circ = \emptyset]$, mais l'implication inverse pas valable en général
\[
P_x \mathrel{|{\circ}|} S_{a,b} \Longleftrightarrow x\le a \text{ ou } y\ge b \Longrightarrow |P_x|^\circ\cap|S_{a,b}|^\circ = \emptyset
\]
mais l'implication inverse pas valable en général
\note{sic $y\ge b$, pour $x\ge b$.}
\[
S_{a,b} \mathrel{|{\circ}|} S_{a',b'} \Longleftrightarrow a<b\le a'<b' \text{ ou } a'<b'\le a<b
\]
\[
\Longrightarrow |S_{a,b}|^\circ \cap |S_{a',b'}|^\circ = \emptyset .
\]
mais implication inverse pas valable en général
\marginal{NB $X \not\lessgtr Y \Longleftrightarrow X = Y$ ou $X \mathrel{|{\circ}|} Y$ ; \ill{} très spécial \ill{} \uncertain{dim} 1 …}

Cet exemple qu'on vérifie directement que $\le$, $\overset{\circ}{\ll}$ sont des relations d'ordre, et $|{\circ}|$ est sym.\ antiréflexive (Mag 1).

Il faut vérifier Mag 2, 3, 4.

Mag 2 OK, on trouve de plus
\[
X \ll Y \Longleftrightarrow
\begin{cases}
X = P_x,\, Y = P_y, & x = y\\
X = P_x,\, Y = S_{a,b}, & a\le x\le b \quad \text{\struck{$x\in[a,b]$ ($=|S_{a,b}|$)}}\\
X = S_{a,b},\, Y = S_{a',b'} : & a'\le a<b\le b' ,
\end{cases}
\]
donc
\[
\begin{cases}
X \ll Y \Longleftrightarrow |X|\subset|Y|\\
\text{si } X,Y\in\mathcal{M}_1, \quad X\ll Y \Longleftrightarrow X \overset{\circ}{\ll} Y
\end{cases}
\]
\[
\begin{aligned}
&\text{Mag 3 :} && X \mathrel{|{\circ}|} Y,\; X'\overset{\circ}{\ll}X,\; Y'\overset{\circ}{\ll}Y \Longrightarrow X'\mathrel{|{\circ}|}Y' && \text{OK}
\end{aligned}
\]
Mag 4 : $X$ \ill{} $Y_0 \Longrightarrow X \mathrel{|{\circ}|} Y$ OK
\note{la relation de Mag 4 est noircie par une surcharge ; le $Y$ porte un petit indice rond.}

Ici encore, les $P_x$ sont \struck{\uncertain{les}} minimaux : à la fois pour $\le$ et pour $\overset{\circ}{\ll}$. Ce sont les seuls éléments minimaux pour $\le$, \struck{\ill{}} mais comme minimaux pour $\overset{\circ}{\ll}$, il y a en plus les éléments minimaux pour $\overset{\circ}{\ll}$, il y a en plus les $S_{a,b}$ tels que $|S_{a,b}|^\circ = \;]a,b[\; = \emptyset$. \uncertain{Toutefois}, les lieux sont les $P_x$, \uncertain{i.e. l'ens.} des lieux s'identifie à $L$.

\page{92}
\note{p. 84 de l'auteur.}

L'axiome de divisibilité \struck{équiv} Mag L 1. \uncertain{équivaut} ici au fait que \struck{\ill{}} l'ens.\ ordonné soit divisible
\[
\text{Mag L 1} \Longleftrightarrow \forall S_{a,b},\; |S_{a,b}|^\circ \overset{\text{déf}}{=} \;]a,b[\; \neq \emptyset
\]
i.e. si $a<b$, $\exists x$ avec $a<x<b$.

Cela implique-t-il Mag L 2 ? \struck{Non, exemple}

Mag L 2 signifie ici deux choses
\[
\begin{aligned}
&\text{a) pour } P_x, S_{a,b} : && x \parallel \;]a,b[\; \Longrightarrow x\le a \text{ ou } x\ge b\\
&\text{b) pour } S_{a,b}, S_{a',b'} : && ]a,b[\; \parallel \;]a',b'[\; \Longrightarrow b\le a' \text{ ou } b'\le a
\end{aligned}
\]
— mais b) vraisemblablement \ill{} \ill{} (i) $L$ soit div.
\note{après a), une ligne barrée « $\Rightarrow$ \ill{} $x\parallel$ \ill{} » ; après b), un renvoi surchargé et entouré, partiellement lisible (« \ill{} $b'$[, \ill{} \ill{} que »), d'où part une flèche vers la condition (i).}

\struck{Dans les deux cas, il faut pour}

\textbf{Lemme} \struck{\ill{}} Pour que Mag L 2 soit valable, il f.\ et il s.\ que \add{(i)} \add{(ii)} que $\forall x\in L$, et $a<b$, on ait
\[
x \parallel \;]a,b[\; \Longrightarrow \text{\struck{$x\parallel\{a,b\}$}}\; x\le a \text{ ou } x\ge b .
\]
\struck{(où pour \uncertain{deux parties} une partie $A$ de $L$, on pose}
\[
\text{\struck{$x\parallel A \Longleftrightarrow \forall y\in A,\; x\parallel y$ (i.e. $x<y$ ou $y<x$))}}
\]
\struck{ou encore, que si $a\in L$ n'est pas un plus grand élément,} \add{si $a\in L$}
\[
L_{\le a} = \text{\struck{$\bigcap_{b\in L_{>a}}$}}
\]
\note{ce dernier encadré, barré d'un trait oblique, est laissé inachevé.}
On encore, que

(iii) a) $\forall a\in L$, t.q.\ $a$ ne soit pas un plus grand élément, on ait
\[
\bigcap_{c>a} L_{\le c} = L_{\le a}
\]
b) $\forall b\in L$ tel que $b$ ne soit pas un plus petit élément, on ait
\[
\bigcap_{c<b} L_{\ge c} = L_{\ge b}
\]

Ce n'est le plus souvent pas satisfait pour les \struck{magasins} \uncertain{ens.\ ordonnés} \struck{\ill{}} associés \struck{\ill{}} à des \uncertain{espaces} topo-logiques, exemple
\drawing{une région allongée en fuseau, hachurée, d'extrémités $0$ et $1$ ; deux bandes verticales la traversent, marquées $a$ et $b$ ; une flèche désigne un point $x$ dans la bande $b$.}

\page{93}
\note{p. 85 de l'auteur.}

Quand le magasin est-il fidèle, i.e.\ quand a-t-on la condition de fidélité
\[
(*) \qquad \mathrm{Omb}^\circ(X) \cap \mathrm{Omb}^\circ(Y) = \emptyset \Longrightarrow X \mathrel{|{\circ}|} Y \quad ?
\]
Il faut déjà que sur $L$
\[
P_x \neq P_y \Longrightarrow x \mathrel{|{\circ}|} y
\]
i.e.\ que $L$ soit \textit{totalement} ordonné. Cela équivaut à $(*)$ pour $X = P_x$, $Y = P_y$. Dans le cas $X = P_x$, $Y = S_{a,b}$, cela signifie $x\notin\;]a,b[\; \Longrightarrow x \mathrel{|{\circ}|} S_{ab}$ i.e.\ $x\le a$ ou $x\ge b$, OK. Dans le cas où $X = S_{a,b}$, $Y = S_{a',b'}$, cela signifie que l'on aura \struck{$S_{ab}$} si on n'a pas $b\le a'$ ou $b'\le a$, i.e.\ on a à la fois \emph{$b>a'$}, \emph{$b'>a$}, alors il y a un \ill{} à $L$ \uncertain{suffisamment} intérieur commun. Je distingue quatre cas, suivant la position relative de $a$ par rapport à $a'$, et de $b$ par r.\ à $b'$
\[
\begin{aligned}
&a'\le a<b\le b' && \text{on prend } S_{ab} \text{ comme \uncertain{\ill{}} intérieur commun}\\
&a\le a'<b\le b' \text{ ou } a'\le a<b'\le b && \text{on prend } S_{a'b} \text{ resp. } S_{ab'}\\
&a\le a'<b'\le b && \text{on prend } S_{a',b'}
\end{aligned}
\]

Donc \uncertain{\emph{lemme}} $\mathcal{M}$ est fidèle ssi $L$ est totalement ordonné.

\textbf{Corollaire} $\mathcal{M}$ est \uncertain{parfait} ssi $L$ est \uncertain{totalement} ordonné et divisible (et alors $\mathcal{M}$ est \ill{} ensembliste).

\page{94}
\note{p. 86 de l'auteur.}

Passons aux questions de support. Dans le cas général, on ne peut se contenter des supports intérieurs sur $L$, faute \uncertain{d'avoir} Mag L 2. \struck{\ill{}} D'abord, étudions les figures (finies).

\struck{$X \not\lessgtr Y$ : si $X = P_x$, $Y = P_y$ : $x\neq y$, \ill{} i.e.\ $\{x,y\}\in\mathrm{Drap}(L)$ ; si $X = P_x$, $Y = S_{a,b}$ : \ill{} i.e.\ $x\notin[a,b]$ : $x\parallel S_{ab}$, $x\in\{a,b\}$ ou $x\parallel$}
\note{ce début de tableau est barré de plusieurs traits obliques.}

Examinons la condition de compatibilité $X\not\lessgtr Y$. Si $X$ (disons ou $Y$) est de la forme $P_x$, donc un lieu
\[
P_x \not\lessgtr Y \Longleftrightarrow P_x = Y \text{ ou } P_x \mathrel{|{\circ}|} Y
\]
(vrai pour des lieux dans un magasin) \struck{donc $\{x,y\}\in\mathrm{Dr}$} donc
\[
\begin{aligned}
P_x \not\lessgtr P_y &\Longleftrightarrow \{x,y\}\in\mathrm{Drap}(L)\\
P_x \not\lessgtr S_{a,b} &\Longleftrightarrow P_x \mathrel{|{\circ}|} S_{a,b} \Longleftrightarrow x\le a \text{ ou } x\ge b
\end{aligned}
\]
D'autre part, $S_{ab}\not\lessgtr S_{a',b'}$ signifie que, ou bien $S_{ab}\parallel S_{a',b'}$, ou $S_{ab} = S_{a',b'}$, \uncertain{ou} $\widetilde{S}_{ab} \cap \widetilde{S}_{a',b'}$ contenu dans $\partial S \cap \partial S'_{a',b'}$ et non vide \struck{et $S_{ab} \mathrel{|{\circ}|} S_{a',b'}$}. Dans ce dernier cas, je dis que l'intersection \add{$\partial S\cap\partial S'$} \struck{est} se réduit à un \add{seul} \struck{point} élément — s'il y en avait deux, on aurait en effet $\partial S = \partial S'$ d'où $S = S'$, car pas \ill{} \ill{} \ill{} cycles. On trouve ainsi
\[
S_{a,b}\not\lessgtr S_{a',b'} \Longleftrightarrow
\begin{cases}
\text{ou (i) } S_{a,b} = S_{a',b'} \quad \text{i.e. } a = a',\, b = b'\\
\text{ou (ii) } S_{ab} \mathrel{|{\circ}|} S_{a',b'}\\
\text{ou (iii) } a<b = a'<b' \text{ ou } a'<b' = a<b
\end{cases}
\]
\note{en (ii), le signe est surchargé et noirci ; à droite, « i.e.\ $a<b<a'<b'$ ou $a'<b'<a<b$ » est barré. Sous (iii), « i.e.\ $S_{ab} \mathrel{|{\circ}|} S_{a',b'}$ » et « mais dans \ill{} $S_{a,b} \mathrel{|{\circ}|} S_{a',b'}$ » sont barrés.}

ou bien
\[
S_{a,b}\not\lessgtr S_{a',b'} \Longleftrightarrow S_{a,b} = S_{a',b'} \text{ ou } S_{a,b} \mathrel{|{\circ}|} S_{a',b'}
\]
\struck{i.e.\ ou} ce dernier cas s'écrivant
\[
a<b\le a'<b' \text{ ou } a'<b'\le a<b
\]

\page{95}
\note{p. 87 de l'auteur.}

Posons
\[
\delta X =
\begin{cases}
\mathbf{x} & \text{si } X = P_x\\
|\partial X| = \{a,b\} & \text{si } X = S_{a,b}
\end{cases}
\]
donc en tous cas
\[
\delta X \subset L, \quad \operatorname{card}\delta X \in \{1,2\}
\]
\struck{Soient} Une condition \textit{nécessaire} pour $X\not\lessgtr Y$ est
\[
X\not\lessgtr Y \Longrightarrow \delta(X)\cup\delta(Y) \in \mathrm{Drap}(L)
\]
\note{$\mathbf{x}$ : le $x$ est écrit en gras, surchargé.}
\marginal{ens.\ des parties finies tot.\ ordonnées de $L$.}

\struck{cette} un drapeau ayant cardinal égal à 1 (si $X = Y\in\mathcal{M}_0$), 2, 3, ou 4 (si $X,Y\in\mathcal{M}_1$ et $X\parallel Y$). Le cas de cardinal 2 est celui de $P_x\not\lessgtr P_y$ avec $x\neq y$, ou de $P_x\not\lessgtr S_{a,b}$ avec $P_x \mathrel{\Delta} S_{ab}$ i.e.\ $x\in\{a,b\}$, ou $X = Y \in \mathcal{M}_1$ ; le cas de cardinal 3 est \uncertain{celui} de $X\parallel Y$ (\uncertain{avec} $P_x, S_{a,b}$ $X\in\mathcal{M}_0$, $Y\in\mathcal{M}_1$ ou inversement, avec $x<a$ ou $x>b$), ou celui de $X,Y\in\mathcal{M}_1$, \uncertain{\ill{}} comme $S_{a,b}$, $S_{b,c}$ avec $a<b<c$ (segments \textit{contigus}). Ainsi trois possibilités
\[
\begin{cases}
\text{lieux égaux ou \struck{\ill{}} disjoints}\\
\text{lieu \struck{\ill{}} incident à un segment, ou disjoint des segments}\\
\text{segments disjoints, segments \textit{adjacents}}
\end{cases}
\]

Mais de façon générale, si on a $T\in\mathrm{Drap}(L)$,
\[
T = \{t_1,t_2,\ldots,t_n\} \quad \text{avec } t_1<t_2<\ldots<t_n ,
\]
définissons les \textit{segments interstitiels} \add{de $T$} (comme les segments $S_{t_1,t_2}$, $S_{t_2,t_3}$, …, $S_{t_{n-1},t_n}$. \struck{Ceux} \add{Les} \struck{dits \ill{}} $P_{t_i}$ sont appelés les \textit{lieux} interstitiels. Ceci dit, on trouve ceci :
\marginal{NB Si $\operatorname{card} T = 1$, \uncertain{segments} \ill{} \ill{} \uncertain{lieux}}

\page{96}
\note{p. 88 de l'auteur.}

\textbf{Proposition} Soit $F\subset\mathcal{M}$ une partie finie \add{non vide} de $\mathcal{M}$, posons \struck{\ill{}}
\[
\delta F = \bigcup_{X\in F} \delta X .
\]
Pour que $F$ soit formée de multistructures mutuellement compatibles, il f.\ et il s.\ que les deux conditions suivantes soient satisfaites
\begin{enumerate}
\item[(i)] $\delta F \in \mathrm{Drap}(L)$ \struck{[NB $\delta F = \emptyset$ ssi $F = \emptyset$]}
\item[(ii)] Les éléments de $F$ sont interstitiels pour $\delta F$.
\end{enumerate}
\marginal{\ill{} \ill{} une partie \ill{} partie \ill{}}

Dém.\ par récurrence sur $\operatorname{card} F$. \struck{\ill{}}

Pour un \struck{\ill{}} \add{$T\in\mathrm{Drap}(L)$} donné, il y a donc une partie \struck{$X\not\lessgtr Y$} de $\mathcal{M}$ maximale, parmi les parties \struck{\ill{}} formées d'él.\ mutuel.\ compatibles et telles que $\delta F = T$ : à savoir
\[
F_T = \text{ensemble des toutes les structures interstitielles de } T .
\]
On l'appellera la figure \uncertain{canonique} définie par $T$.

\struck{Considérons maintenant le cas le plus simple (à part celui où $\operatorname{card} T = 1$) celui de}
\[
\text{\struck{$F = \widetilde{S}_{a,b} = \{S_{a,b}, P_a, P_b\}$ .}}
\]
\struck{On considère les \struck{quatre supports} trois supports}
\[
\text{\struck{$\mathfrak{S}_a = \mathrm{supp}^\circ P_a$, $\mathfrak{S}_b = \mathrm{supp}^\circ P_b$,}}
\]
\struck{$\mathfrak{S}$}\ill{} \struck{$= \mathrm{supp}^\circ S_{a,b}$ .}
\struck{La question des supports concerne donc \ill{} \ill{} de cet ens.\ de trois éléments, \ill{} deux \uncertain{figures}. Il y a \uncertain{manifestement} \ill{} cas distincts non triviaux}
\note{tout ce passage, depuis « Considérons maintenant », est barré de deux longs traits obliques ; sous le premier support, « \ill{} » est raturé.}

\section{Figures d'un magasin associé à un ensemble ordonné}

\page{97}
\note{p. 89 de l'auteur.}

\textbf{2 juillet} Déterminons maintenant \struck{les} les \textit{figures}. Tout drapeau $T$ \add{(\struck{non vide}) de $L$} \add{(détermine une)} figure, $F_T$, savoir \struck{celle \ill{} l'ens}
\[
F_T = \{X\in\mathcal{M} \mid X \text{ interstitiel par } T\} ,
\]
qui se représente graphiquement par
\drawing{un segment horizontal découpé en intervalles par des points marqués $t_1, t_2, t_3, \ldots, t_{n-1}, t_n$.}
(on « oublie » les éléments de $L$ qui ne sont pas dans un des segments $[t_i,t_{i+1}]$, $1\le i\le n-1$). Les cas « dégénérés » sont \struck{$T = \emptyset$, $F_T = \emptyset$}
\[
\begin{aligned}
T = \emptyset &\Longrightarrow F_T = \emptyset && \text{graphisme : ?}\\
T = \{a\} &\Longrightarrow F_T = \{P_a\} && \text{graphisme : un point } a\\
T = \{a,b\} &\Longrightarrow F_T = \widetilde{S}_{a,b} && \text{graphisme : un segment } a \text{---} b
\end{aligned}
\]
puis \uncertain{viennent}
\drawing{deux segments contigus, de sommets $a$, $b$, $c$.}
etc.
\marginal{NB $F_T$ détermine $T$, $T = \delta F_T$}

Sur l'ens.\ des drapeaux \textit{non vides} $\mathrm{Drap}^*(L)$ on met une relation d'ordre
\[
T < T' \quad \text{ssi} \quad \forall t\in T,\; t'\in T', \text{ on a } t<t' ,
\]
qui nous permet de considérer des drapeaux de drapeaux.
\marginal{NB $T<T'$ \ill{} \ill{} $T\parallel T'$ \ill{} $(\mathrm{Drap}^*(L))$ \ill{} \ill{} $T\cap T' = \emptyset$, $T\cup T'\in\mathrm{Drap}^*(L)$}

\textbf{Lemme} Soient $T, T'\in\mathrm{Drap}^*(L)$, il f.\ et il s.\ pour que
\[
F_T \parallel F_{T'} \quad (\Longleftrightarrow F_T \mathrel{|{\circ}|} F_{T'})
\]
que
\[
\underset{\text{dans } \mathrm{Drap}^*(L)}{T\parallel T'} \quad \text{i.e. } \{T,T'\}\in\mathrm{Drap}^*(\mathrm{Drap}^*(L)) .
\]

\page{98}
\note{p. 90 de l'auteur.}

\struck{Si donc on}

\textbf{Corollaire} Soit $\mathcal{T}$ \struck{$\subset$} $\mathrm{Drap}^*(L)$ un ens.\ de drapeaux non vides. Pour que les figures $F_T$ ($T\in\mathcal{T}$) soient deux à deux disjointes, il f.\ et il s.\ que $\mathcal{T}$ soit un drapeau de l'ens.\ ordonné $\mathrm{Drap}^*(L)$.

Dans ce cas, on a donc une figure somme
\[
F_{\mathcal{T}} \overset{\text{déf}}{=} \bigcup_{T\in\mathcal{T}} F_T
\]
on retrouve $F_T$ comme cas particulier
\[
F_T = F_{\{T\}}
\]
\marginal{Les $F_T$}
La figure vide correspond au cas $\mathcal{T} = \emptyset$.

\textbf{Théorème} L'application
\[
\begin{gathered}
\mathcal{T} \longmapsto F_{\mathcal{T}}\\
\mathrm{Drap}(\mathrm{Drap}^*(L)) \longrightarrow \mathrm{Fig\,fin}(\mathcal{M}) = \mathfrak{F}
\end{gathered}
\]
est bijective.

Indication de la démonstration. \struck{Pour} Pour une figure donnée, on récupère les \uncertain{sous}-figures $F_T$ comme les « composantes connexes » (\add{\uncertain{dans un sens}} combinatoire évident) des $F$. Il y a reste à voir que les composantes connexes d'une figure sont bien de la forme $F_T$, puis utiliser le lemme précédent.

\page{99}
\note{p. 91 de l'auteur.}

\drawing{une suite de petites chaînes de segments, de points isolés et de segments, séparées les unes des autres sur une même droite.}

\uncertain{l'}exemple de figure — les composantes connexes « se suivent à la queue-leu-leu ».

Il faudrait déterminer \struck{les \uncertain{supports}} \struck{\ill{}} les supports d'une figure. Pour y arriver, on aura à utiliser

\textbf{Hypothèse H)} \add{a) \ill{} él.\ \uncertain{non} minimal de $L$} \struck{Pour tout él.\ $a$ de $L$, si $a$ n'est pas un plus petit élément, \ill{} \ill{} un plus grand \ill{}}. \struck{\ill{}} ou encore : si $L$ n'a pas de plus petit élément, alors pour \struck{tout} $a\in L$, existe $x\in L$ tel que $x<a$.

b) hypothèse duale.

On aura besoin de cette hypothèse dès le lemme 1

\struck{\textbf{Lemme 1} Soient $X, Y\in\mathcal{M}$ tels que $X\parallel Y$, et soit $\Phi\subset\mathrm{Omb}\,X$, $\Psi\subset\mathrm{Omb}\,Y$. Alors}
\[
\text{\struck{$\mathrm{Supp}^\circ(\Phi\cup\Psi) = \mathrm{Supp}^\circ(\Phi)$}}
\]
\note{ce premier énoncé du lemme 1 est barré de plusieurs traits obliques.}

\textbf{Lemme 1} Soit $X\in\mathcal{M}$. Alors
\[
\mathrm{supp}(X) \Bigl(\overset{\text{déf}}{=} \mathrm{supp}^\circ(\widetilde{X})\Bigr) = \mathrm{Omb}(X) = \{Z\in\mathcal{M} \mid |Z|\subset|X|\}
\]
L'inclusion
\[
\mathrm{Omb}(X) \subset \mathrm{supp}(X)
\]
est triviale, c'est l'autre inclusion qui l'est un peu moins.
\marginal{On \ill{} \ill{} $\mathrm{supp}^\circ(X) = \mathrm{Omb}^\circ(X)$ \ill{} \ill{}, $\mathrm{cosupp}^\circ(X) \overset{\text{déf}}{=} \{Y\in\mathcal{M} \mid$ \ill{} $\}$}
\marginal{Autre présentation : a) Se donner une figure $F$ (\uncertain{revient} à) \ill{} un drapeau $\mathfrak{D}$ ($= \delta F = \bigcup_{T\in\mathcal{T}} T = F_0$) \ill{} b) une partie $F_1$ de l'ens.\ $\mathrm{Int}(\mathfrak{D})$ (ensemble des segments interstitiels de $\mathfrak{D}$ \ill{}) — les segments de $F$ \ill{} \ill{} de la figure ; \uncertain{quand} $\operatorname{card}\mathfrak{D} = n$, il y aura $2^{n-1}$ \ill{} \ill{} qui lui correspondent.}
\marginal{NB Donc \ill{} donné avec \ill{} $2^{n-1}$ figures \ill{}}
\note{les notes de la marge gauche, écrites en long et en partie raturées, ne se lisent que par fragments.}

\page{100}
\note{p. 92 de l'auteur.}

\textbf{Cor 1} \add{$\mathrm{supp} = \mathrm{supp}^\circ(\{P_a\})$}
\[
\mathrm{supp}(\{P_a\}) = \{P_a\}
\]

\textbf{Cor. 2}
\[
\mathrm{supp}(\{S_{a,b}\}) = \Bigl\{ Z \Bigm|
\begin{array}{l}
Z = P_x \text{ avec } a\le x\le b\\
\text{ou } Z = S_{x,y} \text{ avec } a\le x<y\le b
\end{array}
\Bigr\}
\]

\textbf{Cor. 3} \add{$\mathrm{Omb}^\circ(S_{ab}) =$}
\[
\mathrm{supp}^\circ\{S_{ab}\} = \mathrm{supp}(\{S_{ab}\}) \setminus \partial S_{ab} \qquad (\partial S_{ab} = \{P_a, P_b\})
\]
\[
= \Bigl\{ Z \Bigm|
\begin{array}{l}
Z = P_x \text{ avec } a<x<b\\
\text{ou } Z = S_{x,y} \text{ avec } a\le x<y\le b
\end{array}
\Bigr\}
\;\; \text{\struck{$= \mathrm{Omb}^\circ(S_{ab})$}}
\]

On trouve de \uncertain{même}

\textbf{Cor. 4}
\[
\begin{aligned}
\mathrm{supp}^\circ(\{S_{a,b}, P_a\}) &= \mathrm{Omb}^\circ(S_{ab}) \cup \{P_a\}\\
\mathrm{supp}^\circ(\{S_{ab}, P_b\}) &= \mathrm{Omb}^\circ(S_{ab}) \cup \{P_b\}
\end{aligned}
\]

\textbf{Démonstration des Cor 1 et 2.} (équivalents au lemme 1).

\struck{On a, par définition}
\[
\text{\struck{$\mathrm{cosupp}^\circ(P_a) = L_{<a}\cup L_{>a} = \{x\in L \mid x<a \text{ ou } x>a\}$}}
\]
\[
\text{\struck{$\mathrm{supp}^\circ(P_a) = \mathrm{supp}(P_a) = \mathrm{cosupp}^\circ(\mathrm{cosupp}^\circ(P_a))$}}
\]
\struck{Soit $X\in\mathrm{supp}^\circ(P_a)$. On a donc}
\[
\text{\struck{$X \mathrel{|{\circ}|} x \quad \forall x\in\mathrm{cosupp}^\circ(P_a)$}}
\]
\struck{Si $a$ n'est pas plus grand élément}
\note{ce passage est barré de plusieurs traits obliques.}

\uncertain{\emph{Cor. 1}} Soit $X \in \mathrm{supp}^\circ(P_a)$ $(= \mathrm{supp}(P_a))$, i.e.
\[
X \mathrel{|{\circ}|} Y, \quad \forall Y\in\mathrm{cosupp}^\circ(P_a) .
\]
Si $a$ est un plus petit ou plus grand élément, alors
\[
\mathrm{cosupp}^\circ(P_a) = \mathcal{M} - \{P_a\} \quad (\text{car tout } X\in\mathcal{M}
\]
\note{la phrase se poursuit sur la page suivante, au-delà de ce lot.}

\end{document}
